\section{A metric keeps only the properties we use}\label{ch11:sec:metric}
Let $X$ be a set. We call its members \emph{points}, even when they are numbers or vectors. A distance on $X$ should assign a real number $d(x,y)$ to each pair of points $x,y\in X$. In the language of Section~\ref{ch03:sec:functions}, $d$ is a function from $X\times X$ to $\RR$: it takes a pair $(x,y)$ as its input and returns a real number. The following definition lists the three properties that such a function must have before we call it a distance.
\par\noindent\begin{minipage}{\linewidth}
\begin{definition}[Metric]\label{def:metric}
A \emph{metric} on a set $X$ assigns a real number $d(x,y)$ to every pair of points $x,y\in X$ and satisfies these requirements for all $x,y,z\in X$:
\begin{enumerate}
\item $d(x,y)\ge0$, and $d(x,y)=0$ holds exactly when $x=y$.
\item $d(x,y)=d(y,x)$.
\item $d(x,z)\le d(x,y)+d(y,z)$.
\end{enumerate}
A set together with a chosen metric is called a \emph{metric space} and is written $(X,d)$. The number $d(x,y)$ is called the \emph{distance} between $x$ and $y$.
\end{definition}
\end{minipage}\par

The first requirement says that distances are never negative, that every point is at distance zero from itself, and that two different points are always a positive distance apart. The second, called \emph{symmetry}, says that the distance from $x$ to $y$ equals the distance from $y$ to $x$. The third is the triangle inequality of Section~\ref{ch11:sec:distance}, now demanded of $d$: going from $x$ to $z$ by way of a third point $y$ is never shorter than going directly.

Like every definition, this one can be used in the two directions described in Section~\ref{ch01:sec:definition}. If we are told that $d$ is a metric, we may \emph{unpack} the definition and use any of the three requirements in an argument. If we want to show that a new formula is a metric, we must \emph{recognize} it as one, which means checking all three requirements. Calling a formula a distance does not make it one, as the next question shows.

\pauseq{On the real line, let $d(x,y)=(x-y)^2$. Is $d$ a metric?}{No. The first two requirements hold: a square is never negative, $(x-y)^2=0$ exactly when $x=y$, and $(y-x)^2=(x-y)^2$. The triangle inequality fails. Take $x=0$, $y=1$ and $z=2$. Then $d(x,z)=4$, while $d(x,y)+d(y,z)=1+1=2$, so the route through $y$ comes out \emph{shorter} than the direct route. One counterexample is enough to show that $d$ is not a metric.}

\subsection*{Examples}
The real line with $d(x,y)=|x-y|$ is a metric space. The absolute value is never negative, and $|x-y|=0$ exactly when $x-y=0$, that is, when $x=y$. Symmetry, $|x-y|=|y-x|$, and the triangle inequality $|x-z|\le|x-y|+|y-z|$ were both proved in Section~\ref{ch11:sec:distance}. When we treat the real line as a metric space without naming a metric, this is the metric we mean.

The set $\RR^n$ of vectors $x=(x_1,\ldots,x_n)$ from Section~\ref{ch10:sec:vectors} carries several useful metrics. The first is the \emph{Euclidean distance} $d_2(x,y)=\norm{x-y}_2$, the length of the difference vector. In the plane, where $n=2$, it is the ordinary straight-line distance; from $(0,0)$ to $(3,4)$ it is $\sqrt{3^2+4^2}=5$. Its first requirement holds because a length is never negative and is zero only for the zero vector, and $x-y$ is the zero vector exactly when $x=y$. Symmetry holds because $y-x$ has the coordinates of $x-y$ with every sign changed, and changing signs does not change the squares under the root. The triangle inequality was proved in Chapter~\ref{ch:linear}, right after the Cauchy--Schwarz inequality (Theorem~\ref{thm:cs}; see also Exercise~\ref{ex:10.5}). There we showed $\norm{u+v}_2\le\norm{u}_2+\norm{v}_2$ for all vectors $u$ and $v$, and taking $u=x-y$ and $v=y-z$, so that $u+v=x-z$, gave $\norm{x-z}_2\le\norm{x-y}_2+\norm{y-z}_2$. In our new notation this reads $d_2(x,z)\le d_2(x,y)+d_2(y,z)$.

Another useful choice is
\[
d_\infty(x,y)=\max_{1\le j\le n}|x_j-y_j|,
\]
the largest of the $n$ numbers $|x_1-y_1|,\ldots,|x_n-y_n|$. The symbol $\infty$ in the subscript is just part of the traditional name. This metric records the worst error in any single coordinate. For $x=(1,5)$ and $y=(4,3)$, the coordinate differences have sizes $3$ and $2$, so $d_\infty(x,y)=3$, while $d_2(x,y)=\sqrt{9+4}=\sqrt{13}\approx3.61$.

Here is the proof of the triangle inequality for $d_\infty$. Let $x,y,z\in\RR^n$ and fix a coordinate $j$. The triangle inequality for real numbers gives the first step below. For the second step, note that $|x_j-y_j|$ is one of the numbers whose largest is $d_\infty(x,y)$, so $|x_j-y_j|\le d_\infty(x,y)$; in the same way, $|y_j-z_j|\le d_\infty(y,z)$. Thus
\[
|x_j-z_j|\le|x_j-y_j|+|y_j-z_j|\le d_\infty(x,y)+d_\infty(y,z).
\]
The right side is the same number for every $j$. So it is at least as large as each of the numbers $|x_1-z_1|,\ldots,|x_n-z_n|$, and in particular at least as large as the largest of them, which is $d_\infty(x,z)$. This proves $d_\infty(x,z)\le d_\infty(x,y)+d_\infty(y,z)$. Exercise~\ref{ex:11.3} checks the other two requirements.

Which metric to use depends on the question. If a measuring instrument must get \emph{every} coordinate right to within $0.01$, the condition $d_\infty(x,y)\le0.01$ says exactly that. A bound on $d_2(x,y)$ would instead control all the coordinate errors together, through the sum of their squares.

\subsection*{Norms, and spaces without addition}
Both metrics on $\RR^n$ are built from a way of measuring the size of a single vector. For $d_2$ the size is the Euclidean length, since $d_2(x,y)=\norm{x-y}_2$. For $d_\infty$ it is the largest coordinate size $\norm{x}_\infty=\max_j|x_j|$, since $d_\infty(x,y)=\norm{x-y}_\infty$. Measures of size like these are called norms. Precisely, a \emph{norm} on $\RR^n$ assigns to each vector $x$ a real number $\norm{x}$ such that, for all vectors $x,y$ and every real number $c$,
\begin{enumerate}
\item $\norm{x}\ge0$, with equality exactly when $x$ is the zero vector;
\item $\norm{cx}=|c|\,\norm{x}$;
\item $\norm{x+y}\le\norm{x}+\norm{y}$.
\end{enumerate}
The Euclidean length $\norm{x}_2$ is a norm. Its first property was noted in Chapter~\ref{ch:linear}, the second holds because $\sqrt{c^2x_1^2+\cdots+c^2x_n^2}=\sqrt{c^2}\,\norm{x}_2=|c|\,\norm{x}_2$, and the third is the triangle inequality proved there. The largest coordinate size $\norm{x}_\infty$ is also a norm, as you can check coordinate by coordinate in the same way. In the other direction, every norm produces a metric by the rule $d(x,y)=\norm{x-y}$. The first requirement of a metric comes from the first property of the norm. Symmetry comes from the second property with $c=-1$, because $y-x=(-1)(x-y)$. The triangle inequality comes from the third property, because $x-z=(x-y)+(y-z)$.

A metric, however, does not need points to be added or subtracted at all. It needs only one number $d(x,y)$ for each pair of points. For this reason the rest of the chapter works in a general metric space. The arguments then apply also to sets in which addition makes no sense, or in which a sum can fall outside the set.

The most common way to obtain such a set is to take part of a larger metric space. If $(X,d)$ is a metric space and $C$ is a subset of $X$, then $d$, used only for pairs of points of $C$, is a metric on $C$. (In the language of Section~\ref{ch03:sec:injsurj}, this is the restriction of $d$ to $C\times C$.) Indeed, the three requirements hold for all points of $X$, so in particular they hold for all points of $C$. For example, the interval $[0,1]$ with distance $|x-y|$ is a metric space. Its points can be added as real numbers, but a sum can leave the interval, as $0.7+0.6=1.3$ does. Arguments that use only distances never meet this difficulty.
